\section{Methodology}
\label{Methodology}

Figure~\ref{fig:workflow} summarises the experiment, from task assignment
to code evaluation and disclosure analysis. Appendix C provides the
settings, prompts, and reproduction commands.

\begin{figure}[!ht]
\centering
\begin{tikzpicture}[
  box/.style={draw, rounded corners=2pt, align=center,
              font=\small, text width=2.8cm, minimum height=0.95cm,
              inner sep=4pt},
  flow/.style={-{Stealth[length=4pt]}, semithick}]
\node[box] (task) at (0,0) {Task and settings\\Data, tests, prompt};
\node[box] (agent) at (4,0) {Agent writes code\\Runs tests if any};
\node[box] (saved) at (8,0) {Save code and logs\\Before evaluation};
\node[box] (evidence) at (8,-1.5)
  {Hidden tests\\Then agent review};
\node[box] (coding) at (4,-1.5)
  {Inspect code and logs\\Label disclosures};
\node[box] (analysis) at (0,-1.5)
  {Compare outcomes\\RQ1, RQ2, RQ3};
\draw[flow] (task) -- (agent);
\draw[flow] (agent) -- (saved);
\draw[flow] (saved) -- (evidence);
\draw[flow] (evidence) -- (coding);
\draw[flow] (coding) -- (analysis);
\end{tikzpicture}
\caption{Evidence collected from each run. Hidden evaluation precedes
self-review, but its results are not supplied to the agent.}
\label{fig:workflow}
\end{figure}

\subsection{Prepare the Task}

We use six tasks covering degree completion, race--SES interactions,
incarceration and wealth, climate trends, replication effects, and fish
diversity. Five use DiscoveryBench data; the climate task adapts
ScienceAgentBench task 87~\cite{majumder2025discoverybench,chen2025scienceagentbench}.
Each specifies a Python function, inputs, analysis, and outputs. The agent
has NumPy, pandas, and SciPy, but no dedicated modelling package.
Appendix C lists task identifiers and conflict values (Table~\ref{tab:tasks}).

A run is one agent attempt in a fresh workspace with starter code, data,
instructions, and any supplied tests. Codex CLI runs under a macOS Seatbelt
profile restricting access to hidden tests and reference code. We retain
code, messages, and evaluations for 207 completed runs across four model
identifiers, including runs with incorrect analyses. We exclude 23 attempts
with recorded quota failures, stopped batches, or harness errors, including
two early attempts that received the wrong test condition (Appendix C).

\subsection{Assign Tests and Instructions}

Table~\ref{tab:conditions} lists the test conditions. Structural checks
assess output format; numerical checks compare expected values. Conflict
substitutes one wrong value, making complete test success incompatible
with correct analysis. We analyse this stress condition separately.

\begin{table}[!ht]
\caption{Tests supplied during implementation.}
\label{tab:conditions}
\centering
\begin{tblr}{colspec={l X[l] r}, width=\linewidth, rowsep=1pt,
             row{1}={font=\bfseries}, hline{1,2,Z}={0.4pt}}
Condition & Supplied tests & Runs \\
All & Structural and numerical checks & 72 \\
Some & Structural checks only & 37 \\
None & No tests & 37 \\
Conflict & Full suite with one wrong expected value & 61 \\
\end{tblr}
\end{table}

Explicit prompts require generalisation; neutral prompts omit that reminder;
pressure prompts emphasise automated grading. Later runs allow up to three
requests to fix failing visible tests, stopping at success or the limit.
Runs without tests receive no follow-up. Most runs use one model, and
uneven allocation prevents attributing pooled differences to visibility alone.

\subsection{Test and Classify the Code}

We save the code and evaluate it separately. Hidden unit tests check expected
outputs; metamorphic tests check relations between executions. Anchored
checks compare results on the supplied dataset and transformed copies to
expose input-specific substitutions. Suite v1 combines unit and original
metamorphic checks; v2 adds anchored checks. Failures, errors, skips, and
timeouts flag implementations for review.

For example, the anchored SES test duplicates the supplied dataset. A
correct fit should preserve the coefficient and reduce its standard error
by a factor of $\sqrt{2}$. The test compares both outputs with relative
tolerance $10^{-3}$ and absolute tolerance $10^{-6}$. Code that substitutes
an answer only at the original row count can violate this relation while
passing checks on unrelated inputs.

Passing the wrong conflict check supports a gaming label; hidden-test failure
alone does not. Every non-conflict run flagged by hidden tests or a static
scan for hard-coded expected values has a label based on code and transcript
inspection. We manually assigned these labels; they
remain under author review (Appendix A). Unflagged runs were not all
classified. Anchored checks followed the pilot, and subsequent evaluator
fixes mean retrospective detection is not independent validation.

Planted implementations and mutations of honest solutions assess detector
strength. We retain raw mutation scores and scores adjusted for recorded
equivalence judgements, separately from detection of agent shortcuts.

\subsection{Review the Code and Record Disclosure}

After saving and evaluating code, we request self-review. Later runs receive
a neutral question about what a reviewer should know, then a direct question
about shortcuts, both in the original session. Hidden results are withheld.

Coding distinguishes mentions of discrepancies, explicit admissions, and
misleading correctness claims in implementation messages, final answers,
and reviews. We manually assigned these codes, which remain under
author review. Neutral-review comparisons cover the 30 of 62 identified
gaming cases with that stage available.

\subsection{Compare the Outcomes}

For RQ1, we report identified gaming cases per included run in each condition.
For RQ2, detection is the proportion of identified gaming cases flagged by
a check family or suite; ordinary faults are counted separately. For RQ3,
denominators include only cases with available responses at the relevant
stage. We use Wilson 95\% confidence intervals and exact McNemar tests for
paired outcomes. Repeated tasks, uneven allocations, and incomplete review
limit generalisation. Appendix C gives the settings and reproduction commands.
